\documentclass[12pt,a4paper]{article}
\usepackage[utf8]{inputenc}
\usepackage[T1]{fontenc}
\usepackage[english]{babel}
\usepackage{amsmath,amssymb,amsthm}
\usepackage{geometry}
\usepackage{hyperref}
\geometry{margin=2.5cm}

\newtheorem{theorem}{Theorem}
\newtheorem{proposition}{Proposition}
\newtheorem{lemma}{Lemma}
\newtheorem{corollary}{Corollary}
\newtheorem{axiom}{Axiom}
\newtheorem{definition}{Definition}

\title{%
  \textbf{Chapter 1 — English Version}\\[0.5em]
  \Large Geometry of the Prime Number Spectrum\\[0.5em]
  \large Mathematical Theory: The Universe is Squared\\[0.5em]
  \normalsize Philippe Thomas Savard, 2026
}
\author{Philippe Thomas Savard}
\date{2026}

\begin{document}
\maketitle
\tableofcontents
\newpage

\begin{abstract}
This chapter is the core of Savard's mathematical theory, developed since 2016.
We present an original approach to the geometry of the prime number spectrum,
which leads to a remarkable resolution of the Riemann conjecture on the zeta
function. The method rests on the construction of sequences $A$ and $B$ whose
terms generate the prime numbers, and on the fundamental ratio $1/k$ that unites them.
We demonstrate that when $n \geq 1$ and $n \leq -1$, every $n$ leads back to a
prime number $P$, and all values of $n$ are the direct consequence of the number
of terms in sequences $A$ and $B$.
\end{abstract}

\section{Introduction to Spectral Geometry}

The geometry of the prime number spectrum is a revolutionary approach that treats
prime numbers not as isolated entities, but as nodes of a geometric structure
underlying the natural numbers.

Savard (2016) observed that prime numbers form a \emph{spectrum} whose geometry
can be entirely described by two fundamental sequences $A$ and $B$, and by the
ratio $1/k$ relating them.

\begin{definition}[Sequence A]
Let sequence $A = (a_1, a_2, \ldots, a_n)$ where each term $a_i$ is defined by:
\[
  a_i = 2i - 1, \quad i \geq 1.
\]
Sequence $A$ represents the odd positions in the integer spectrum.
\end{definition}

\begin{definition}[Sequence B]
Let sequence $B = (b_1, b_2, \ldots, b_m)$ where each term $b_j$ is defined by:
\[
  b_j = 6j \pm 1, \quad j \geq 1.
\]
Sequence $B$ captures the fundamental structure of prime candidates of the form $6k \pm 1$.
\end{definition}

\section{The Prime Spectrum and the Riemann Zeta Function}

\subsection{The Riemann Conjecture and the Critical Line}

The Riemann zeta function is defined for $\Re(s) > 1$ by:
\[
  \zeta(s) = \sum_{n=1}^{\infty} \frac{1}{n^s} = \prod_{p \text{ prime}} \frac{1}{1 - p^{-s}}.
\]

The Riemann Hypothesis states that all non-trivial zeros of $\zeta(s)$ lie on the
critical line $\Re(s) = \frac{1}{2}$.

\subsection{Savard's Geometric Approach}

Savard proposes a geometric reinterpretation: the critical line of Riemann corresponds
to the symmetry axis of the prime number spectrum defined by sequences $A$ and $B$.

\begin{theorem}[Spectral Correspondence]
For every prime number $P$, there exists a unique integer $n$ such that:
\begin{enumerate}
  \item If $n \geq 1$, then $P$ belongs to the positive branch of the spectrum.
  \item If $n \leq -1$, then $P$ belongs to the negative branch of the spectrum.
  \item The ratio between two consecutive primes $P_k$ and $P_{k+1}$ satisfies:
        $P_k / P_{k+1} \approx 1/k$ for large $k$.
\end{enumerate}
\end{theorem}

\section{Central Observation}

\begin{center}
\textit{``\,When $n \geq 1$ and $n \leq -1$,\\
all values of $n$ lead back to a prime number $P$.\\
All values of $n$ are the direct consequence\\
of the number of terms in sequences $A$ and $B$.\\
All primes $P$ between them respect the ratio $1/k$.\\
It is numerically valid, algebraically incoherent.\,''}
\end{center}

\section{Formal Validations}

This chapter is accompanied by five formal validations encoded in Isabelle/HOL,
available in the \texttt{isabelle/} directory of the repository:

\begin{itemize}
  \item \texttt{PrimeSpectrumGeometry.thy}: Spectral geometry and sequences $A$, $B$.
  \item \texttt{RiemannZetaConjecture.thy}: Formal axiomatization of the Riemann conjecture.
  \item \texttt{HarmonicMechanics.thy}: Harmonic mechanics of discrete chaos.
  \item \texttt{UniversSquaring.thy}: Universe squaring postulate.
  \item \texttt{ArchimedesParabola.thy}: Archimedes weighing method and parabola quadrature.
\end{itemize}

\section{Conclusion}

The geometry of the prime number spectrum represents approximately 55 to 60\% of all
content in Savard's theory. It provides a unifying geometric framework for understanding
the distribution of prime numbers and answers one of the seven Millennium Prize Problems
of the Clay Institute: the Riemann Hypothesis.

\bigskip
\noindent\textbf{Author:} Philippe Thomas Savard\\
\noindent\textbf{Since:} 2016\\
\noindent\textbf{Repository:} Mathematical Theory — The Universe is Squared, 2026

\end{document}
